\subsection{Evolución del almacenamiento de credenciales}
La versión inicial almacenaba contraseñas en texto plano. Se reemplazó este mecanismo por Argon2id en el servidor: el registro calcula el hash antes de persistirlo y el inicio de sesión verifica la contraseña contra la cadena almacenada. Se eligió una configuración de 19 MiB de memoria, dos iteraciones y paralelismo uno, de acuerdo con la recomendación mínima de OWASP~\cite{owasp-passwords}. La biblioteca genera un salt aleatorio independiente de 16 bytes por operación y codifica los parámetros, el salt y el resultado en la misma cadena~\cite{node-argon2}. La estructura de las siete tablas se conserva.

Las cuentas existentes se actualizan con una migración transaccional que modifica únicamente el valor de la contraseña. Los hashes Argon2id ya presentes no se vuelven a transformar. Si falla una actualización, se revierten los cambios de la transacción. El procedimiento preserva los identificadores y las contraseñas originales de acceso; no se implementó una comparación alternativa de texto plano en el login. La ejecución debe repetirse en cada instalación existente.

Se mantienen el bloqueo temporal, los roles y la comprobación de cuenta activa. Las respuestas públicas no incluyen contraseña ni hash. Las pruebas verifican salts diferentes para contraseñas iguales, rechazo de credenciales incorrectas, uso de la contraseña completa, repetición de la migración y reversión ante un fallo.

\subsection{Sesión web y mejora propuesta}
Las páginas utilizan el JWT de la API y lo conservan durante la sesión de la pestaña. El cierre de sesión elimina la copia local; no revoca individualmente un JWT ya emitido. El almacenamiento del token en JavaScript conserva un riesgo frente a XSS, por lo que se restringen los scripts al mismo origen y los datos de usuarios se incorporan como texto. Un despliegue público requiere HTTPS.

Se propone WebAuthn como segundo factor para administradores, con soporte en la documentación de OWASP sobre autenticación multifactor~\cite{owasp-mfa}. El hash protege el almacenamiento, pero no impide el uso de una contraseña robada. Un segundo factor agrega una comprobación al acceso administrativo. La mejora no se implementó porque requiere definir el alta, la recuperación y las pruebas del segundo factor; permanece como una evolución posterior.
